\chapter{Esercizi di riepilogo}


\begin{exercise}

Supponiamo di avere uno strumento di misura in cui un valore analogico, soggetto
ad un rumore elettronico gaussiano con media nulla e deviazione standard $\sigma$
nota precisamente a priori, viene convertito in un valore digitale tramite un
convertitore analogico-digitale (ADC) a $12$ bit che, per definizione, mostra sul display
un numero intero compreso tra $0$ e $2^{12}-1 = 4095$.

Si discura l'errore di misura da assegnare alla lettura digitale dello strumento
al variare di $\sigma$ (espressa in unit\`a di conteggi di ADC). In quali condizioni
ha senso fare la media di misure ripetute per minimizzare l'errore di misura?


% A digital balance displays measurements rounded to the nearest gram. A student therefore models a displayed value (m_0) as uniformly distributed over

% [
% [m_0-0.5\ \mathrm g,;m_0+0.5\ \mathrm g].
% ]

% The instrument, however, has internal Gaussian electronic noise with standard deviation (0.35\ \mathrm g), applied before rounding.

% A reference object has true mass (m), and the displayed value is

% [
% D=\operatorname{round}(m+\epsilon),
% \qquad
% \epsilon\sim\mathcal N(0,0.35^2\ \mathrm g^2).
% ]

%     Derive (P(D=k\mid m)).
%     Suppose (D=58\ \mathrm g). Assuming a locally uniform prior for (m), derive the posterior density of (m).
%     Is the posterior uniform between (57.5) and (58.5\ \mathrm g)?
%     Calculate or numerically estimate its posterior variance.
%     Explain precisely when the conventional (1/\sqrt{12}\ \mathrm g) uncertainty is justified and when it is not.

% Subtlety: a finite display resolution does not by itself imply a uniform conditional distribution for the true value.

\end{exercise}